\centerline{\bf \Huge II Отбор на подготовительные курсы 7-IT}

\begin{flushright}
        \scriptsize{\textbf{\textit{Именно математика даёт надёжнейшие правила: тому кто им следует — тому не опасен обман чувств. \\Леонард Эйлер}}}
        
\end{flushright}

\problem{\textbf{1}} Начальник записал на доске следующую сумму равную $X$: 

$2010\cdot2009 + 2008\cdot2007 + 2006\cdot2005 + 2004\cdot2003 + 2002\cdot2001 + 2000\cdot1999= X$

Чему равна последняя цифра числа $X$?

\hspace{3mm}

\problem{\textbf{2}} В отборе участвуют только шестиклассники, рассаженные в шесть кабинетов. В первом кабинете шестиклассников в два раза меньше, чем во втором, во втором на четыре шестиклассника меньше, чем в третьем, в третьем столько же, сколько и в четвертом, в четвертом на 50\% меньше, чем в пятом, а в пятом в четыре раза больше чем в шестом. Может ли в пятом кабинете сидеть шестиклассников ровно 33\% от общего числа шестиклассников, участвующих в отборе?

\hspace{3mm}

\problem{\textbf{3}} Чему равно следующее выражение?

$50\cdot\left(\dfrac{(1,238+2,762) \cdot 0,1}{(36,487-34,237): 2,8125}+\dfrac{(4,36-1,16) \cdot 0,3125}{0,2 \cdot(47,8-45,55): 0,225}\right)\cdot$ 

$\cdot\left(39,86-\dfrac{8 \dfrac{4}{5}}{6,934+38 \dfrac{13}{40}: 2 \dfrac{1}{2} \cdot 0,2}-17\right)$ = ?

\hspace{3mm}

\problem{\textbf{4}} Аркадий и Савр решили устроить лёгкую пробежку вокруг ЭТЛ по часовой стрелке. Каждый раз, когда Аркадий обгоняет Савра, он называет вслух число, равное количеству обгонов им Савра (сразу после первого обгона называет число 1, сразу после второго обгона называет число 2 и так далее). Чему равна сумма всех чисел, названных Аркадием, если известно, что мальчики начали одновременно с одного места и Аркадий быстрее Савра в три раза?

\hspace{3mm}

\centerline{\textbf{\textit{Задания 5,6,7 на обратной стороне}}}

\newpage
\hspace{3mm}

\problem{\textbf{5}} Среди 2026 учеников ЭлМШ есть \textit{шиншиллы}, которые всегда говорят правду, и \textit{лососи}, которые всегда лгут (есть хотя бы одна \textit{шиншилла} и хотя бы один \textit{лосось}). Какое наибольшее количество \textit{лососей} может быть, если каждый сказал: "Всего среди 2026 учеников не менее 11 \textit{лососей}"?

Есть ли логическая ошибка в условии задания? Если есть, объясните её и предложите решение этой проблемы. Если нет, решите задание.

\hspace{3mm}

\problem{\textbf{6}} Однажды Елена решила отправиться из Оргакина (посёлок в Ики-Бурульском районе) в разностороннее путешествие. На своё усмотрение за одно действие она может поехать либо строго на юг, либо строго на север. За первое действие Елена проехала 1км, а за каждое последующее действие в четыре раза больше предыдущего (4км, 16км, 64км, 256км и так далее). Может ли Елена вернуться в Оргакин за некоторое количество таких действий (Оргакин не объёмный и является точкой на плоскости)?

\hspace{3mm}

\problem{7} Адьян пришел на III ежегодную Математическую Олимпиаду ЭлМШ x ЭТЛ, на которой было предложено 5 задач. Адьян очень торопился и мог потратить на Олимпиаду максимум час. Задачи, придуманные Расулом Владимировичем, Адьян читает за 2 минуты, а решает за 35 минут. Задачи, придуманные Анджуром Аркадьевичем, он читает 2 минуты, а решает за 15 минут. Задачи, придуманные Очиром Владимировичем, он читает 10 минут и решает тоже за 10 минут. Очир Владимирович предложил две задачи, Анджур Аркадьевич одну, а Расул Владимирович две. Какое количество баллов Адьян мог набрать, если задачи Очира Владимировича он мог отличить, даже не читая, а отличить задачи Анджура Аркадьевича от задач Расула Владимировича может только прочитав их, при этом, читая, сначала ему попадаются задачи только Расула Владимировича. За задачи Очира Владимировича ученики получали по 4 балла, Анджура Аркадьевича - по 5 баллов, Расула Владимировича - по 8 баллов.

\hspace{3mm}

\centerline{\textbf{\textit{На каждое задание нужно дать развёрнутое решение, а не просто ответ!}}}